% !TeX root = main.tex

\section{Conclusion}
\label{sec:conclusion_generale}

Ce projet avait un objectif principal : lire automatiquement des archives foncières très diverses, extraire les informations utiles et les verser sur Géofoncier, le tout sans aucun transfert de données vers l'extérieur du cabinet.


\subsection{Ce que le projet a produit}

Avant ce PFE, retrouver un dossier ancien demandait de consulter physiquement les registres, de chercher dans des boîtes ou d'appeler un autre site du cabinet. Verser un document sur Géofoncier était une démarche ponctuelle, manuelle, réalisée seulement quand le temps le permettait. L'outil produit ici change ces deux réalités : la recherche d'antériorités peut désormais s'effectuer en quelques secondes depuis un navigateur, et le versement sur Géofoncier s'intègre directement dans le flux de validation, sans étape supplémentaire pour l'opérateur.

Ce résultat a été obtenu sans recours à un service cloud et sans budget matériel spécifique : les modèles tournent sur les postes existants du cabinet, les données ne quittent jamais le réseau local. La contrainte de confidentialité imposée par le secret professionnel a orienté chaque choix technique, et elle a été tenue.


\subsection{Ce qui reste à faire}

Ce projet a surtout démontré que l'approche fonctionne. Sur les types de documents les plus courants du cabinet, la chaîne produit des fiches exploitables, le versement sur Géofoncier est opérationnel, et les opérateurs ont pris l'outil en main rapidement lors des sessions de test. C'est le résultat principal du PFE.

La suite dépend de la mise en pratique. Le volume réel des archives est bien plus grand que le jeu de test : traiter les 23\,600 dossiers du siège d'Aubenas demandera plusieurs mois de travail régulier par les opérateurs. C'est cet usage continu, plus que tout ajustement technique, qui déterminera la valeur finale de l'outil pour le cabinet. Les pistes d'amélioration identifiées au chapitre précédent, comme le modèle HTR pour le Géomètre B, la boucle d'apprentissage actif, la migration de l'interface, restent des idées ouvertes, réalisable indépendamment selon les priorités du cabinet.


\subsection{Bilan personnel}

Ce projet a confirmé que la problématique traitée ici n'est pas propre au cabinet GEO-SIAPP. Un grand nombre de documents sont stockés dans les archives des cabinets de géomètres, et la plupart des équipes les recherchent manuellement. Le besoin existe à grande échelle, et l'approche développée ici, même incomplète, montre qu'il est possible d'y répondre avec des outils accessibles, sans infrastructure cloud et sans budget matériel important.

Des bases en Python existaient avant ce projet, mais la distance entre écrire des scripts simples et assembler une arborescence de traitement d'images fonctionnel sur des documents dégradés est réelle. La difficulté principale n'était pas technique, c'était de partir d'une problématique large, de la cadrer, et de produire quelque chose utile au cabinet en six mois. Les outils d'IA générative ont été une aide pour avancer plus vite sur les parties de code les moins familières, notamment le traitement d'image avec OpenCV et l'intégration des modèles via Ollama, avant tout application sur des exemples concrets pour ne pas diffuser ces mêmes plans.

Ce projet a également rappelé que l'IA n'est pas une solution en soi. C'est un ensemble d'outils qui, bien choisis et bien assemblés, peuvent répondre à des problèmes très concrets. Mais chaque modèle a ses limites, et ces limites apparaissent vite dès qu'on sort des conditions idéales. Le travail de fiabilisation et de contrôle, décrit au chapitre~\ref{sec:fiabilisation}, est au moins aussi important que le choix des modèles.


\subsection{Ouverture}

Ce projet s'inscrit dans un mouvement plus large. Le 48\textsuperscript{e} Congrès des géomètres-experts, qui s'est tenu à Nantes du 23 au 25 juin 2026, a été entièrement consacré à l'impact de l'intelligence artificielle sur les technologies et les expertises de la profession. Plus de 1\,400 professionnels s'y sont réunis, ce qui en fait le plus grand rassemblement de l'Ordre depuis sa création. À l'issue de cette édition, une charte d'utilisation de l'IA, votée le 19 mai 2026, est entrée en application le 1\textsuperscript{er} septembre 2026. Elle repose sur quatre principes : transparence vis-à-vis du client, protection des données, traçabilité, et maintien de la responsabilité du professionnel sur le livrable final.

Ce cadre correspond à ce que ce projet a cherché à respecter : l'outil produit des propositions, l'opérateur valide. Aucun versement ne se fait sans relecture humaine. La question de la place de l'IA dans le métier de géomètre-expert ne se limite d'ailleurs pas aux archives. Le traitement de nuages de points, la génération de procès-verbaux, l'analyse d'images issues de drones : autant de domaines où des approches similaires pourraient être utiles. Ce PFE est, à sa modeste échelle, un cas d'usage concret parmi ceux que la profession commence à explorer.



